%! Polar Function Graphs — Unit 3, Lesson 3.14 — Slides (Beamer)
\documentclass[aspectratio=169, 11pt]{beamer}
\usepackage{precalculus-beamer}   % provides \navyheader{} and \sectionlabel[color]{}

\begin{document}

% TITLE SLIDE
{
\setbeamercolor{background canvas}{bg=navy}
\begin{frame}[plain]
  \vspace{2.8cm}\hspace{0.55cm}%
  \begin{minipage}{0.9\textwidth}
    {\color{goldacc}\bfseries\small LESSON 3.14}\\[0.5cm]
    {\color{white}\bfseries\LARGE Polar Function Graphs}\\[1.2cm]
    {\color{skymid}\small AP Precalculus~$\cdot$~Unit 3: Trigonometric and Polar Functions}
  \end{minipage}
\end{frame}
}

% FRAME 1 — What is a polar function?
\begin{frame}[t, plain]
  \navyheader{What Is a Polar Function?}

  \begin{block}{Definition}
    A \textbf{polar function} is $r = f(\theta)$ where:
    \begin{itemize}
      \item $\theta$ is the \textbf{input} --- the angle measure from the positive $x$-axis
      \item $r$ is the \textbf{output} --- the radius (distance from the origin)
    \end{itemize}
  \end{block}

  \vspace{0.4cm}
  \textbf{Each $(\theta, r)$ pair locates a point:}
  \begin{itemize}
    \item Start at the origin, face the direction $\theta$ from the positive $x$-axis
    \item Move $r$ units in that direction (or $|r|$ units in the opposite direction if $r < 0$)
  \end{itemize}

  \vspace{0.4cm}
  \begin{block}{Example}
    For $r = 2\sin\theta$ at $\theta = \pi/2$: \quad $r = 2\sin(\pi/2) = 2(1) = 2$\\
    The point is 2 units from the origin along the ray at $\theta = \pi/2$ (straight up).
  \end{block}
\end{frame}

% FRAME 2 — Reading from a table
\begin{frame}[t, plain]
  \navyheader{Reading a Polar Function from a Table}

  \textbf{Table of values for $r = 2\sin\theta$:}

  \vspace{0.3cm}
  \renewcommand{\arraystretch}{1.4}
  \begin{tabular}{|c|c|c|c|c|c|c|c|}
  \hline
  $\theta$ & $0$ & $\pi/6$ & $\pi/3$ & $\pi/2$ & $2\pi/3$ & $5\pi/6$ & $\pi$ \\
  \hline
  $r$ & $0$ & $1$ & $1.73$ & $2$ & $1.73$ & $1$ & $0$ \\
  \hline
  \end{tabular}

  \vspace{0.4cm}
  \begin{itemize}
    \item As $\theta$ goes $0 \to \pi/2$: $r$ \textbf{increases} from 0 to 2
    \item As $\theta$ goes $\pi/2 \to \pi$: $r$ \textbf{decreases} from 2 to 0
    \item At $\theta = 0$ and $\theta = \pi$: $r = 0$ --- the curve \textbf{passes through the pole}
    \item Maximum radius: $r = 2$ at $\theta = \pi/2$
  \end{itemize}

  \vspace{0.3cm}
  \begin{block}{The result}
    These input-output pairs trace a \textbf{circle} of radius 1 centered at $(0, 1)$ in
    the polar plane!
  \end{block}
\end{frame}

% FRAME 3 — Connecting rectangular and polar representations
\begin{frame}[t, plain]
  \navyheader{Connecting Rectangular and Polar Graphs}

  \textbf{The rectangular graph of $r = f(\theta)$ is a road map for the polar curve:}

  \vspace{0.3cm}
  \begin{itemize}
    \item \textbf{Zeros} on the rectangular graph ($r = 0$) $\Rightarrow$ the polar curve
      \textbf{passes through the pole}
    \item \textbf{Maxima} ($r$ at its largest positive value) $\Rightarrow$
      the \textbf{tip of a petal} (farthest from origin)
    \item \textbf{Negative} $r$ $\Rightarrow$ the point plots in the
      \textbf{opposite direction} from $\theta$
  \end{itemize}

  \vspace{0.4cm}
  \begin{block}{Rose curve: $r = \sin(2\theta)$}
    The rectangular graph completes \textbf{4 zero-to-max-to-zero arcs} in $[0, 2\pi]$
    (2 positive, 2 negative). Since $n = 2$ is even, the polar curve has $2n = 4$ petals.
  \end{block}

  \vspace{0.3cm}
  \textbf{General rule:} $r = \sin(n\theta)$ has $n$ petals (odd $n$) or $2n$ petals (even $n$).
\end{frame}

% FRAME 4 — Domain restriction
\begin{frame}[t, plain]
  \navyheader{Domain Restriction}

  \begin{block}{Key Idea (EK 3.14.A.2)}
    The domain of $\theta$ can be restricted to trace only a \textbf{portion} of
    the polar curve. Smaller interval $=$ fewer petals (or part of a curve).
  \end{block}

  \vspace{0.4cm}
  \textbf{Example with $r = \sin(2\theta)$:}

  \vspace{0.3cm}
  \begin{itemize}
    \item $\theta \in \left[0,\, \dfrac{\pi}{2}\right]$:
      one petal traced (in the first quadrant)
    \item $\theta \in [0,\, \pi]$:
      two petals traced
    \item $\theta \in [0,\, 2\pi]$:
      all four petals traced
  \end{itemize}

  \vspace{0.4cm}
  \begin{block}{How to read domain restriction}
    Find the sub-interval on the rectangular graph. Count how many complete
    zero-to-max/min-to-zero arcs fit. Each arc $=$ one petal.
  \end{block}
\end{frame}

\end{document}
